\documentclass[11pt,a4paper]{article}

\def\courseTitle{Industrial Security (Bachelor)}
\def\sectionNumber{02}
\def\sectionTitle{ICS and SCADA Fundamentals}
\def\sheetKind{Solution Sheet}

% =====================================================================
% Shared preamble for exercise sheets and solution sheets.
%
% Define BEFORE \input{}-ing this file:
%   \def\courseTitle{...}
%   \def\sectionNumber{...}
%   \def\sectionTitle{...}
%   \def\sheetKind{Exercise Sheet}   or  Solution Sheet
%
% Optional:
%   \def\courseAuthor{...}
%
% The caller must issue \documentclass{...} (typically article) BEFORE
% loading this preamble.
% =====================================================================

\usepackage[utf8]{inputenc}
\usepackage[T1]{fontenc}
\usepackage[english]{babel}
\usepackage[a4paper, margin=2.2cm]{geometry}
\usepackage{graphicx}
\usepackage{listings}
\usepackage{xcolor}
\usepackage{tikz}
\usepackage{booktabs}
\usepackage{array}
\usepackage{enumitem}
\usepackage{titlesec}
\usepackage{fancyhdr}
\usepackage{hyperref}
\usepackage{amsmath}
\usepackage{amssymb}
\usepackage{tcolorbox}
\tcbuselibrary{skins, breakable}
\usepackage{needspace}

% Discourage solo lines split across pages.
\widowpenalty=10000
\clubpenalty=10000
\raggedbottom

% =====================================================================
% Shared TikZ styles for the Industrial Security lecture series.
% Loaded by every slide deck and exercise sheet that uses TikZ.
% =====================================================================

\usetikzlibrary{
  arrows.meta,
  shapes,
  shapes.geometric,
  shapes.symbols,
  shapes.multipart,
  positioning,
  fit,
  calc,
  backgrounds,
  decorations.pathmorphing,
  decorations.pathreplacing,
  decorations.text,
  patterns,
  matrix,
  shadows
}

% --- colour palette --------------------------------------------------
\definecolor{isOT}{HTML}{1F4E79}     % OT (deep blue)
\definecolor{isIT}{HTML}{6F2DA8}     % IT (purple)
\definecolor{isSafety}{HTML}{C00000} % safety red
\definecolor{isNet}{HTML}{2E7D32}    % network green
\definecolor{isAttack}{HTML}{B23A48} % attack red
\definecolor{isAccent}{HTML}{E08E0B} % amber
\definecolor{isMute}{HTML}{6B6B6B}   % grey
\definecolor{isLight}{HTML}{EDF2F8}  % light background
\definecolor{isPLC}{HTML}{2C5282}
\definecolor{isHMI}{HTML}{2F855A}
\definecolor{isSCADA}{HTML}{6B46C1}

% --- generic node styles --------------------------------------------
\tikzset{
  isnode/.style={
    draw, rounded corners=2pt, minimum height=8mm, minimum width=18mm,
    font=\small, align=center, thick
  },
  isbox/.style={isnode, fill=isLight},
  plc/.style={isnode, fill=isPLC!15, draw=isPLC, thick},
  hmi/.style={isnode, fill=isHMI!15, draw=isHMI, thick},
  scada/.style={isnode, fill=isSCADA!15, draw=isSCADA, thick},
  sensor/.style={isnode, fill=isNet!10, draw=isNet, minimum height=6mm, minimum width=14mm, font=\scriptsize},
  actuator/.style={isnode, fill=isAccent!15, draw=isAccent, minimum height=6mm, minimum width=14mm, font=\scriptsize},
  firewall/.style={isnode, fill=isAttack!15, draw=isAttack, thick},
  server/.style={isnode, fill=isMute!15, draw=isMute!80, thick},
  switch/.style={isnode, fill=isLight, draw=black!60, minimum height=6mm, minimum width=14mm, font=\scriptsize},
  workstation/.style={isnode, fill=white, draw=isOT, thick},
  attacker/.style={isnode, fill=isAttack!20, draw=isAttack, thick, font=\small\bfseries},
  inetnode/.style={draw, ellipse, minimum width=24mm, minimum height=10mm, font=\scriptsize, fill=isLight, thick},
  process/.style={isnode, fill=isOT!15, draw=isOT, thick, minimum width=24mm},
  zone/.style={
    draw, rounded corners=4pt, thick, dashed, inner sep=4mm
  },
  conduit/.style={-{Stealth[length=2.5mm]}, thick, isNet!80!black},
  attack/.style={-{Stealth[length=2.5mm]}, thick, isAttack, dashed},
  flow/.style={-{Stealth[length=2.5mm]}, thick, isOT!70!black},
  netline/.style={thick, black!60},
  axisline/.style={-{Stealth[length=2mm]}, thick, black!70},
  tag/.style={font=\scriptsize\itshape, fill=white, inner sep=1pt},
  level/.style={
    draw, fill=isLight, minimum width=0.85\linewidth, minimum height=8mm,
    font=\small, anchor=west
  },
  brace/.style={decorate, decoration={brace, amplitude=4pt, raise=2pt}},
  legendbox/.style={draw, rounded corners=2pt, fill=isLight, inner sep=2mm, font=\scriptsize, align=left}
}

% --- handy macros ----------------------------------------------------
% \plcicon, \hmiicon etc as simple labelled boxes; useful inline.
\newcommand{\plcicon}[1][PLC]{\tikz[baseline=-0.6ex]{\node[plc, minimum width=10mm, minimum height=5mm, font=\scriptsize, inner sep=1pt]{#1};}}
\newcommand{\hmiicon}[1][HMI]{\tikz[baseline=-0.6ex]{\node[hmi, minimum width=10mm, minimum height=5mm, font=\scriptsize, inner sep=1pt]{#1};}}
\newcommand{\fwicon}[1][FW]{\tikz[baseline=-0.6ex]{\node[firewall, minimum width=10mm, minimum height=5mm, font=\scriptsize, inner sep=1pt]{#1};}}


\providecommand{\courseAuthor}{Industrial Security Lecture Series}
\providecommand{\courseInstitute}{Department of Computer Science}
\providecommand{\companyLogo}{logo}

% --- Corporate branding overrides ------------------------------------
% Picks up logo / author / brand colours from the repo-root branding.tex
% when present; falls back to the defaults above otherwise.
\IfFileExists{../../../branding.tex}{% =====================================================================
% Corporate / institutional branding -- single file to customise the
% look of the entire course (slides, notes, exercises, labs).
%
% HOW TO REBRAND IN UNDER A MINUTE:
%   1. Replace `branding/logo.pdf` with your own logo
%      (PDF preferred; PNG and JPG also work -- name it `logo.<ext>`).
%   2. (Optional) change the two colours below to your brand palette.
%   3. (Optional) override the author/institute lines below.
%   4. Re-run `make` at the repo root. That's it.
%
% Everything in this file has sensible defaults; you can delete a line
% to fall back to the built-in defaults, or comment it out.
%
% This file is loaded automatically by the slides, exercise, and lab
% preambles -- you never need to \input it manually.
% =====================================================================

% --- Logo --------------------------------------------------------------
% Path is resolved from each leaf-build directory; the default points at
% the shared `branding/` folder at the repo root. If you put your logo
% somewhere else, set the full or relative path here (without extension):
\renewcommand{\companyLogo}{../../../branding/logo}

% --- Author / institute (appears on the title page footer) ------------
\renewcommand{\courseAuthor}{}
\renewcommand{\courseInstitute}{}

% --- Brand colours ----------------------------------------------------
% slideAccent      = primary brand colour (title bands, accent rule)
% slideSecondary   = secondary brand colour (section badge, highlight)
% Both use HTML hex codes. Keep enough contrast against white text.
\definecolor{slideAccent}{HTML}{00205B}      % deep navy
\definecolor{slideSecondary}{HTML}{00A0DC}   % light blue accent
}{}

% --- listings -------------------------------------------------------
\definecolor{codebg}{rgb}{0.96,0.96,0.96}
\lstset{
  backgroundcolor=\color{codebg},
  basicstyle=\ttfamily\footnotesize,
  breaklines=true,
  frame=single,
  numbers=left,
  numberstyle=\tiny\color{black!50},
  showstringspaces=false,
  tabsize=2
}

% --- header / footer ------------------------------------------------
\pagestyle{fancy}
\fancyhf{}
\renewcommand{\headrulewidth}{0.4pt}
\renewcommand{\footrulewidth}{0pt}
\lhead{\small \courseTitle}
\rhead{\small Section \sectionNumber{} -- \sheetKind}
\cfoot{\thepage}

% --- task / solution boxes -----------------------------------------
% `before={...}` guards every task/solution box with \needspace so the
% box doesn't start with only one or two lines of breathing room at the
% bottom of a page. If less than ~9 lines remain, LaTeX bumps the box
% to the next page; otherwise it starts here and breaks normally if it
% runs long.
% Boxes are `breakable` so very long solutions don't blow the page,
% but `before={\needspace{14\baselineskip}}` pushes them to a fresh
% page whenever fewer than ~14 lines remain. That covers the vast
% majority of single-question splits in practice.
\newtcolorbox{taskbox}[1]{
  enhanced, breakable,
  colback=isLight, colframe=isOT,
  fonttitle=\bfseries\color{white},
  coltitle=white,
  colbacktitle=isOT,
  title={#1},
  arc=2pt, boxrule=0.6pt,
  left=2mm, right=2mm, top=2mm, bottom=2mm,
  before={\par\vspace{2mm}\needspace{14\baselineskip}\par},
  after={\par\vspace{2mm}}
}

\newtcolorbox{solutionbox}{
  enhanced, breakable,
  colback=isNet!5, colframe=isNet,
  fonttitle=\bfseries\color{white},
  coltitle=white,
  colbacktitle=isNet,
  title={Solution},
  arc=2pt, boxrule=0.6pt,
  left=2mm, right=2mm, top=2mm, bottom=2mm,
  before={\par\vspace{2mm}\needspace{14\baselineskip}\par},
  after={\par\vspace{2mm}}
}

\newtcolorbox{hintbox}{
  enhanced, breakable,
  colback=isAccent!10, colframe=isAccent,
  fonttitle=\bfseries\color{white}, coltitle=white, colbacktitle=isAccent,
  title={Hint},
  arc=2pt, boxrule=0.6pt
}

\newtcolorbox{warnbox}{
  enhanced, breakable,
  colback=isAttack!8, colframe=isAttack,
  fonttitle=\bfseries\color{white}, coltitle=white, colbacktitle=isAttack,
  title={Caution},
  arc=2pt, boxrule=0.6pt
}

% --- section formatting --------------------------------------------
\titleformat{\section}{\Large\bfseries\color{isOT}}{\thesection}{0.5em}{}
\titleformat{\subsection}{\large\bfseries\color{isOT!80!black}}{\thesubsection}{0.5em}{}

% --- title block macro ----------------------------------------------
\newcommand{\sheetheader}{%
  \begin{center}
    {\Large\bfseries \courseTitle}\\[0.3em]
    {\large \sheetKind{} -- Section \sectionNumber}\\[0.2em]
    {\large \sectionTitle}\\[0.2em]
    \rule{\linewidth}{0.4pt}
  \end{center}
  \vspace{0.5em}
}


\begin{document}
\sheetheader

\section*{How to use this sheet}
Each box gives a model answer suitable for a TA briefing. Open-ended
exercises (2.4, 2.6, 2.7) include rubric hints in brackets at the end of
the solution -- students do not have to match the model answer word for
word, but the rubric items must be present.

\begin{solutionbox}
\textbf{Exercise 2.1 -- Worst-Case Scan-Cycle Latency.}
The worst case stacks: the sensor edge is missed by the current cycle, so
we wait one full cycle for the next input-read phase, then another cycle
for the program to act on the new value, then the output card delay.
Input filter: 3\,ms. One missed cycle plus the
following execute/update cycle: 12 + 12 = 24\,ms.
Output card: 2\,ms. Total worst-case
latency: 3 + 24 + 2 = 29\,ms. The
budget of 40\,ms is met with 11\,ms of margin. To
shave latency, configure the input event on a hardware interrupt OB
(OB40-family on S7-1500) so the reject does not have to wait for the
cyclic scan -- or use an HSC channel. To make it predictably worse, add a
chatty OPC UA poll or large data block to the cyclic OB so housekeeping
time pushes the worst-case cycle nearer to the watchdog.
\end{solutionbox}

\begin{solutionbox}
\textbf{Exercise 2.2 -- Read This Ladder Rung.}
(1) On Start press, \texttt{X1}=1, \texttt{X2}=1 (Stop NC unpressed),
\texttt{X3}=1 (tank not full): \texttt{Y1} energises. From the next scan
on, the \texttt{Y1} seal-in keeps the rung true even after \texttt{X1}
drops back to 0, so the pump runs continuously.
(2) When the tank reaches the high-level switch, \texttt{X3} goes to 0,
breaking the series path; \texttt{Y1} de-energises and the seal-in
collapses on the same cycle. The pump stops automatically -- this is the
intended high-level shutoff.
(3) Cutting the Stop wire makes \texttt{X2} read as 0 (no current through
the NC contact's input circuit). The NC inversion makes \texttt{]/[}
evaluate to 1, so the rung still works -- which is the \emph{wrong}
failure mode: a Stop button whose wire is cut can no longer stop the
pump. Stop must be wired so a broken wire \emph{de-energises} the output
(fail-safe). Most plants therefore wire Stop as a hard-wired series
contact in the output circuit \emph{outside} the PLC, not as a soft input.
\end{solutionbox}

\begin{solutionbox}
\textbf{Exercise 2.3 -- Why Isolated SIS.}
(1) Two technical reasons:
(a) \emph{Common-cause failure}: any firmware bug, runtime fault, memory
corruption, or malicious logic change in the shared M580 would defeat
both the control \emph{and} the protection function in a single event.
(b) \emph{Diversity and SIL certification}: a SIL~3 demand mode requires
proven hardware reliability figures (PFD$_\text{avg}$) and a certified
development process; a general-purpose BPCS controller is not certified
for that and its firmware change cadence is incompatible with SIS
recertification cycles. A typical SIS uses TMR (triple modular redundant)
voting, which the M580 does not provide.
(2) Regulatory: IEC~61511 (the process-industry application of IEC~61508)
mandates that the SIS be ``separate and independent'' from the BPCS, and
that any shared element be justified and analysed. EN~61511 is the
EU-harmonised version; US sites often cite ISA~84.00.01.
(3) TRITON compromised a Schneider Triconex by writing safety logic over
the proprietary TriStation protocol. With shared BPCS+SIS on one
controller, the attacker would need only one foothold to disable both
control and safety. With a separate Triconex on its own engineering
network, the attacker must additionally pivot into the SIS engineering
path \emph{and} bypass the physical keyswitch -- a much higher bar.
\end{solutionbox}

\begin{solutionbox}
\textbf{Exercise 2.4 -- Windows XP HMI.}
A defensible top-five (any reasonable five; rubric items in brackets):
(a) Put the HMI behind a small dedicated firewall (e.g.\ Fortinet
FortiGate Rugged 60F or Phoenix Contact mGuard) with a strict allow-list:
only the PLC IPs on Modbus/S7 ports outbound, only the operator-VLAN
inbound on RDP from named jump host. (b) Disable every Windows service
that is not strictly needed (Server, Workstation, Browser, Messenger);
remove network printer sharing. (c) Apply application allow-listing with
McAfee Application Control (formerly Solidcore) or Microsoft AppLocker
configured in enforce mode -- XP cannot get patches, so we stop unknown
binaries from executing instead. (d) Remove all USB-mass-storage drivers
or enforce DeviceLock; the engineering laptop uses a one-way file diode
(e.g.\ Owl) or a sanitised USB station. (e) Move the HMI off the office
domain and onto a local workgroup with unique local credentials; back the
HMI image up nightly to an offline target so a re-image is one hour, not
one week.
The cheapest per-euro is the firewall plus allow-list (one rule set
removes 95\,\% of the internet-exposed attack surface). Reject:
``install antivirus'' -- modern AV signatures no longer ship for XP,
performance is poor, and most XP-era exploits live below the AV
hook anyway.
[Rubric: at least one network control, at least one host-hardening
control, at least one backup/recovery control, named products, one
explicit rejection with reason.]
\end{solutionbox}

\begin{solutionbox}
\textbf{Exercise 2.5 -- OPC UA on Pump P-101.}
(1) The engineer needs: routable IP reachability to the sensor on TCP
port 4840 (default OPC UA binary); an OPC UA client such as UAExpert; the
server's endpoint URL; if the server enforces a security policy other
than \texttt{None}, an X.509 client certificate trusted by the server
(typically pushed via the GDS or copied into the server's \texttt{trusted}
folder); a valid user token (anonymous, username/password, or X.509
user).
(2) The attacker needs the same three things -- reachability, a client,
and the endpoint. In a typical out-of-the-box deployment, sensor vendors
ship with \texttt{SecurityPolicy=None} and \texttt{Anonymous} access
enabled by default, so the attacker needs nothing extra beyond a network
sniff to learn the endpoint URL. There is effectively no authentication
gap between the two parties.
(3) Disable the \texttt{None} security policy on the server, require
\texttt{Basic256Sha256} with mutual certificate trust, and create a
non-anonymous user token for the engineer. The engineer's UAExpert still
works once her certificate is approved; the attacker now needs a trusted
client certificate \emph{and} valid credentials.
\end{solutionbox}

\begin{solutionbox}
\textbf{Exercise 2.6 -- Historian for Billing.}
Three controls you would add for the billing historian, but not for a
research-only one:
(a) \emph{Append-only write path with cryptographic auditing}: enable PI
Audit Trail (PI Server feature) plus tamper-evident logging via the
PI~System Auditing module, with logs shipped off-box to a SIEM such as
Splunk. (Lives in the PI Server and the SIEM.)
(b) \emph{Time integrity}: harden the historian's NTP source (use an
internal stratum-1 GPS clock from Meinberg or Microsemi rather than
public pool servers); enable PI's secure-time feature so out-of-order
timestamps are flagged. (Lives in the OS and the PI Server.)
(c) \emph{Read-only, signed export to the billing system}: replace the
PI~OLEDB pull with a one-way PI-to-PI replication into a second
``billing-PI'' instance in a separate security zone, with PI~Web~API
served over mutual TLS to the billing application; the billing system
verifies row hashes recorded in PI~Asset~Framework. (Lives in PI~AF, the
network DMZ, and the billing system.)
For a research historian you do not care: scientists rerun their
analytics, can fix bad timestamps, and no money rides on a single value.
If the regulator can audit unannounced, also add immutable retention of
the raw 1\,s tag values for the entire billing-relevant period
(typically 6+ years), so the auditor can recompute the totals.
[Rubric: each control names where it lives and one concrete product or
feature; the regulator add-on is about retention/non-repudiation, not
just ``add more logging''.]
\end{solutionbox}

\begin{solutionbox}
\textbf{Exercise 2.7 -- Open-Book: S7-1500 Scan Times.}
Model answer (Siemens S7-1500 System Manual, current edition at time of
writing, e.g.\ 12/2023, chapter ``Cycle and reaction times''):
The minimum cycle time configurable for the cyclic main OB
(OB~1) is 1\,ms (entered in the CPU's
``Cycle'' properties in TIA Portal). The default maximum cycle-time
monitoring (watchdog) is 150\,ms; configurable
from 1\,ms up to 6\,000\,ms.
Numbers vary slightly between SKUs (a CPU 1518 has higher throughput than
a 1511) but the configurable range is the same.
Siemens exposes a minimum so the user can guarantee a fixed cycle for
deterministic control loops (the CPU idles to fill the cycle); the
maximum is a watchdog -- if the program ever runs past it, the CPU goes
to STOP rather than silently violating real-time guarantees. Free-running
at maximum speed would make jitter unbounded and break downstream timing.
[Rubric: numbers are in the right order of magnitude (min $\approx$ 1\,ms,
max $\le$ 6\,000\,ms), the citation names a specific section/table and
edition date, and the explanation mentions determinism plus the watchdog
purpose. Exact figures from the most current manual are acceptable; the
student must cite the edition they read.]
\end{solutionbox}

\begin{solutionbox}
\textbf{Exercise 2.8 -- Misplaced Seal-In Bug.}
(1) The seal-in is around \emph{Start itself}, not around the whole
start/stop chain. Critically, the second rung does not use the output
\texttt{Run} as the seal contact -- it uses \texttt{Start} again, which
is the pushbutton input that returns to 0 as soon as the operator lets
go. There is therefore no latch at all: the rung is just
\texttt{(Start~OR~Start)~AND~NOT~Stop} = \texttt{Start~AND~NOT~Stop}.
(2) The operator presses Start: while the button is held, \texttt{Run}
energises and the motor turns. The moment she releases the button,
\texttt{Run} drops out and the motor stops -- behaving like a momentary
pushbutton, not a latched start/stop. Pressing Stop does nothing
additional, because the motor was never latched in the first place.
(3) Corrected rung -- the seal contact must be the output \texttt{Run}
itself, in parallel with \texttt{Start}:
\begin{verbatim}
    |   Start    Stop(NC)               Run  |
    |--] [---+--]/[--------------------(   )-|
    |        |                                |
    |   Run  |                                |
    |--] [---+                                |
\end{verbatim}
Now once \texttt{Run} energises for one scan, its own NO contact holds
the rung true through the seal branch until Stop is pressed.
\end{solutionbox}

\begin{solutionbox}
\textbf{Exercise 2.9 -- Component Mapping.}
(1) \textbf{PLC}. RTU is wrong because the description specifies factory
floor, not remote telemetry; HMI does not drive I/O; SIS would not be
the primary controller, only the protection layer.
(2) \textbf{SIS}. A PLC or DCS could trip a valve but does not provide
the SIL-rated independence required for ``trips independent of the
BPCS''. Engineering WS and Historian have no I/O.
(3) \textbf{Engineering Workstation}. The HMI does not download firmware,
the historian does not hold project files, and the PLC itself does not
hold the credentials -- it accepts them.
(4) \textbf{RTU}. PLCs are not designed for long unreliable links and do
not buffer DNP3 events; SCADA Server sits at the master end, not the
remote.
(5) \textbf{Historian}. A PLC has no long-term storage; an HMI keeps
short-term trends but not ten years; an Engineering WS is not the data
serving system.
(6) \textbf{HMI}. Historians provide raw data but not mimics; the
SCADA Server can drive HMIs but the operator's physical viewing surface
is the HMI; PLCs have no display.
\end{solutionbox}

\begin{solutionbox}
\textbf{Exercise 2.10 -- Lab Notes (OpenPLC + Modbus TCP).}
Expected observations:
\begin{itemize}
  \item OpenPLC Runtime listens on Modbus TCP port 502 by default.
  \item The HMI uses function code \texttt{0x01} or \texttt{0x02} to poll
    discrete inputs/coils (Start, Stop, Run), and \texttt{0x05} (write
    single coil) or \texttt{0x0F} (write multiple coils) to set Start/Stop.
    Holding-register access uses \texttt{0x03} read and \texttt{0x06}~/
    \texttt{0x10} write.
  \item The traffic carries no authentication, no integrity check, no
    encryption -- a Wireshark dissector shows full plaintext register
    addresses and values. Anyone on the VLAN can replay or forge a
    write to coil~0 (Start) and start the conveyor.
  \item The corrected ladder behaves identically to Exercise 2.8: brief
    Start press latches the Run output until Stop is asserted.
\end{itemize}
\end{solutionbox}

\end{document}
